\documentclass[a4paper,11pt]{article}

\usepackage[czech]{babel}
\usepackage[T1]{fontenc}
\usepackage[utf8]{inputenc}
\usepackage{amsmath,amssymb}
\usepackage{geometry}
\usepackage{enumitem}

\geometry{a4paper,margin=17mm}
\pagestyle{empty}
\setlength{\parindent}{0pt}
\setlength{\parskip}{4pt}
\setlist[enumerate]{leftmargin=*,itemsep=4pt,topsep=4pt}

% Sazba vychází ze sablony/testy/A4.tex; jedna varianta na jedné A4.
% Zdroj VŠECH zadání: přímé zadání učitele v konverzaci dne 27. 10. 2026,
% zpráva „Vytvoř čtvrtletní písemný test z matematiky na téma posloupnosti“.
% Jednoznačné identifikace: Varianta A, úlohy 1a–d, 2–5, 6a–b;
% Varianta B, úlohy 1a–d, 2–5, 6a–b. Obsah ani číselné hodnoty neměněny.
% Zařazení: 4_rocnik/tematicke_plany/SPS.pdf, s. 1, body 1.1–1.5,
% září–říjen a 1. písemná práce. V dostupné učebnici samostatná kapitola
% posloupností nebyla nalezena. Související podklad pro značení:
% 4_rocnik/pracovni_listy/pracovniListGeoPosl.tex (zejména úlohy 66 a 71).
% Úlohy pracovního listu nejsou zdrojem převzatých zadání tohoto testu.
% Bodování a řešení jsou výhradně v souboru s příponou _reseni.tex.
\newcommand{\TestTitle}{Aritmetická a geometrická posloupnost}
\newcommand{\TestHeader}[1]{%
    \begin{center}
        {\Large\bfseries ČTVRTLETNÍ PÍSEMNÁ PRÁCE}\\[2mm]
        {\bfseries \TestTitle}\\[1mm]
        {\bfseries Varianta #1}
    \end{center}
    \vspace{2mm}
    Jméno a příjmení: \rule{9cm}{0.4pt}\hfill Třída: \rule{2cm}{0.4pt}\\[3mm]
    Datum: \rule{3cm}{0.4pt}\hfill Počet bodů: \rule{2cm}{0.4pt}\hfill Známka: \rule{2cm}{0.4pt}\\[2mm]
    \hrule\vspace{4mm}
}
\newcommand{\AnswerSpace}[1]{\vspace{#1}}

\begin{document}
\TestHeader{A}

\begin{enumerate}
    \item Rozhodněte, zda jsou následující posloupnosti aritmetické,
    geometrické, nebo ani jedno. U aritmetické určete diferenci $d$,
    u geometrické kvocient $q$.

    \begin{tabular}{@{}p{0.47\linewidth}p{0.47\linewidth}@{}}
        a) $7,\ 12,\ 17,\ 22,\ldots$ & b) $3,\ 9,\ 27,\ 81,\ldots$ \\[5mm]
        c) $2,\ 5,\ 10,\ 17,\ldots$ & d) $16,\ {-8},\ 4,\ {-2},\ldots$
    \end{tabular}
    \AnswerSpace{4mm}

    \item V aritmetické posloupnosti platí $a_4=11$, $a_{13}=47$.
    Určete diferenci $d$, první člen $a_1$, člen $a_{20}$ a součet $S_{20}$.
    \AnswerSpace{19mm}

    \item V geometrické posloupnosti s kladným kvocientem platí
    $a_2=6$, $a_5=162$. Určete kvocient $q$, první člen $a_1$,
    člen $a_7$ a součet $S_6$.
    \AnswerSpace{20mm}

    \item Je dána aritmetická posloupnost $4,\ 9,\ 14,\ldots$
    Určete index $n$, pro který platí $a_n=149$, a následně určete součet $S_n$.
    \AnswerSpace{17mm}

    \item Deset čísel tvoří aritmetickou posloupnost s diferencí $d=2$.
    První, třetí a sedmé číslo této posloupnosti tvoří tři po sobě jdoucí
    členy geometrické posloupnosti. Určete všech deset čísel.
    \AnswerSpace{24mm}

    \item Na začátku roku uložíme do banky $40\,000$ Kč. Vklad je úročen
    $6\,\%$ ročně a úroky se vždy na konci roku připíší k vkladu.
    Daň z úroků neuvažujte.
    \begin{enumerate}[label=\alph*),topsep=1pt,itemsep=1pt]
        \item Určete, kolik peněz bude na účtu po 5 letech.
        \item Určete, po kolika celých letech bude na účtu poprvé více
        než $60\,000$ Kč.
    \end{enumerate}
    \AnswerSpace{20mm}
\end{enumerate}

\newpage
\TestHeader{B}

\begin{enumerate}
    \item Rozhodněte, zda jsou následující posloupnosti aritmetické,
    geometrické, nebo ani jedno. U aritmetické určete diferenci $d$,
    u geometrické kvocient $q$.

    \begin{tabular}{@{}p{0.47\linewidth}p{0.47\linewidth}@{}}
        a) $18,\ 13,\ 8,\ 3,\ldots$ & b) $2,\ {-6},\ 18,\ {-54},\ldots$ \\[5mm]
        c) $1,\ 4,\ 9,\ 16,\ldots$ & d) $81,\ 27,\ 9,\ 3,\ldots$
    \end{tabular}
    \AnswerSpace{4mm}

    \item V aritmetické posloupnosti platí $a_5=-3$, $a_{14}=42$.
    Určete diferenci $d$, první člen $a_1$, člen $a_{20}$ a součet $S_{20}$.
    \AnswerSpace{19mm}

    \item V geometrické posloupnosti s kladným kvocientem platí
    $a_2=12$, $a_5=324$. Určete kvocient $q$, první člen $a_1$,
    člen $a_7$ a součet $S_6$.
    \AnswerSpace{20mm}

    \item Je dána aritmetická posloupnost $91,\ 84,\ 77,\ldots$
    Určete index $n$, pro který platí $a_n=-42$, a následně určete součet $S_n$.
    \AnswerSpace{17mm}

    \item Deset čísel tvoří aritmetickou posloupnost s diferencí $d=4$.
    První, třetí a sedmé číslo této posloupnosti tvoří tři po sobě jdoucí
    členy geometrické posloupnosti. Určete všech deset čísel.
    \AnswerSpace{24mm}

    \item Na začátku roku uložíme do banky $50\,000$ Kč. Vklad je úročen
    $5\,\%$ ročně a úroky se vždy na konci roku připíší k vkladu.
    Daň z úroků neuvažujte.
    \begin{enumerate}[label=\alph*),topsep=1pt,itemsep=1pt]
        \item Určete, kolik peněz bude na účtu po 6 letech.
        \item Určete, po kolika celých letech bude na účtu poprvé více
        než $70\,000$ Kč.
    \end{enumerate}
    \AnswerSpace{20mm}
\end{enumerate}
\end{document}
